\documentclass[11pt]{article}

\usepackage[margin=0.92in]{geometry}
\usepackage[T1]{fontenc}
\usepackage{lmodern}
\usepackage{microtype}
\usepackage{amsmath,amssymb}
\usepackage{booktabs}
\usepackage{tabularx}
\usepackage{array}
\usepackage{longtable}
\usepackage{enumitem}
\usepackage{xcolor}
\usepackage{xurl}
\usepackage{hyperref}

\hypersetup{
  colorlinks=true,
  linkcolor=blue!55!black,
  citecolor=blue!55!black,
  urlcolor=blue!65!black,
  pdftitle={ERGT: Attention-Free Native Geodesic Transport for Controlled Long-Horizon Relational Reasoning},
  pdfauthor={Jalal Jafari and Nasibeh Aramideh}
}

\setlist{nosep,leftmargin=*}
\setlength{\emergencystretch}{2em}
\newcommand{\ERGT}{\textsc{ergt}}
\newcommand{\TF}{\textsc{Transformer}}
\newcommand{\R}{\mathbb{R}}
\newcommand{\firstpagecodeavailability}[1]{%
  \begingroup
  \renewcommand{\thefootnote}{}\footnote{#1}%
  \addtocounter{footnote}{-1}%
  \endgroup
}
\newcolumntype{Y}{>{\raggedright\arraybackslash}X}
\newcolumntype{P}[1]{>{\raggedright\arraybackslash}p{#1}}

\title{\textbf{ERGT: Attention-Free Native Geodesic Transport for\\
Controlled Long-Horizon Relational Reasoning}}

\author{
Jalal Jafari\\
\small Ph.D. in Photonics; Board Member, NovoXpert AI Group\\
\small Artificial Intelligence Research and Development\\
\small \href{mailto:j.jafari2009@gmail.com}{\texttt{j.jafari2009@gmail.com}}
\and
Nasibeh Aramideh, Ph.D.\\
\small Ph.D. in Mathematics; Research Member, NovoXpert AI Group\\
\small Artificial Intelligence Research and Development\\
\small \href{mailto:aramidehnasibeh@gmail.com}{\texttt{aramidehnasibeh@gmail.com}}
}
\date{July 18, 2026}

\begin{document}
\maketitle
\firstpagecodeavailability{\textbf{Code availability.} The complete source
code and executable reproduction package accompanying this article, including
the Colab notebook, are publicly available at
\href{https://github.com/jalaljafari2009/ERGT-paper}{\texttt{github.com/jalaljafari2009/ERGT-paper}}.}

\begin{abstract}
We introduce Executable Relational Geometry and Transport (\ERGT), an
attention-free architecture in which relational computation is organized by
input-conditioned geometry rather than self-attention. \ERGT learns a mapping
from raw tokens to a persistent information field $\Psi$, whose pairwise
observables induce multiple directed world geometries. These geometries then
become the computational substrate for reasoning: GeoTransport carries typed
state along admissible paths, and a geodesic boundary solver accepts the
least-cost candidate consistent with propagation, memory, payload, and terminal
constraints. This native answer path uses neither a Transformer nor a supplied
graph, target compiler, or generic executor.

Across four independent runs, \ERGT converged and satisfied every registered
native stability contract. Beyond its training range, mean accuracy at 32 hops
was $0.955$, compared with $0.500$ for a strengthened direct Transformer given
the same raw input and a closely matched parameter count. Frozen
same-checkpoint interventions caused sharp losses when geodesic action,
propagation admissibility, typed transport, geometric memory, or terminal
measurement was removed. Randomizing or shuffling the geometry, or restricting
inference to one world, also failed to recover the intact result. The evidence
supports a bounded architectural claim: a $\Psi$ field formed from raw tokens
can become an executable geometry
and provide a measurable advantage in controlled long-horizon relational
reasoning.
\end{abstract}

\noindent\textbf{Keywords:} relational reasoning; representation learning;
information geometry; geodesic transport; causal intervention; multiworld
geometry; attention-free neural architecture; long-horizon generalization

\section{Introduction}

Modern sequence models have shown how effective content-dependent token
interaction can be for representation learning. In the Transformer, that
interaction is implemented by self-attention \cite{vaswani2017attention}.
However, long-horizon relational reasoning imposes additional structural
requirements. Entities must remain distinguishable, directed and typed
relations must compose in the correct order, intermediate state must survive
repeated propagation steps, and the model should decline to answer when no
valid support path exists. These requirements can be learned implicitly, but a
generic sequence backbone does not make admissibility, conservation, or
rejection part of the computation itself.

Executable Relational Geometry and Transport (\ERGT) is built around that
distinction. Its central premise is that the arrangement of information in an
input can generate the geometry on which reasoning takes place. Rather than
starting from attention scores or a graph supplied with the task, the model
forms a persistent information field $\Psi$ from the raw tokens. This field
keeps identity and relational evidence tied to the same source state, providing
the substrate from which all later geometry is derived.

Pairwise measurements of $\Psi$ produce several directed geometries, or
\emph{worlds}, each acting as a lens on a different relational tendency. The
term is operational: a world contributes support, direction, and traversal
cost; it does not denote a literal physical universe. GeoTransport propagates
typed state along the resulting admissible paths while preserving entity,
operator, payload, and memory information. A native geodesic solver then finds
the least-cost valid route, and boundary measurement accepts a candidate only
when the transported state satisfies the terminal constraints. The complete
path is
\[
\begin{aligned}
  \text{raw tokens}
  &\longrightarrow \Psi
  \longrightarrow \text{multiworld geometry} \\
  &\longrightarrow \text{GeoTransport}
  \longrightarrow \text{geodesic solution}
  \longrightarrow \text{boundary decision}.
\end{aligned}
\]
No graph or geometry is supplied at the input boundary, and the tested native
path does not delegate its decision to a generic compiler or graph executor.
The relation support, path costs, transported state, and acceptance conditions
must all arise from the field constructed from the same raw sequence.

The design is motivated by a specific long-horizon failure mode. When a
solution depends on many typed relations, repeated global remixing may blur
which state belongs to which entity and which path. \ERGT instead treats
composition as repeated local transport under shared admissibility and
conservation rules. Its inductive bias is not simply a wider context window;
it is the preservation of executable relational paths as reasoning depth
increases.

The experimental evaluation follows the same logic. It first tests whether the
full route from raw-token field formation to terminal measurement can converge
and satisfy the registered native stability contracts. Frozen checkpoints are
then evaluated on freshly generated, raw-disjoint relational chains well
beyond the training horizon. Finally, frozen targeted interventions, geometry
corruption, and single-world restrictions test whether the solved answers
actually depend on the proposed mechanism.

The primary baseline is a strengthened direct Transformer trained on the same
raw sequences and matched target families. Staged training and auxiliary
supervision improve its opportunity to learn the task, while its inference
path remains purely Transformer-based: it receives no \ERGT graph, geometry,
transport state, compiler, executor, or identity side channel. This preserves a
clear comparison between two complete inference strategies rather than lending
the proposed solver to the baseline.

The contributions of this work are:
\begin{enumerate}
  \item a persistent $\Psi$ field, built from raw tokens, that supplies the
  internal identity and pairwise evidence from which relation geometry is
  induced rather than given;
  \item an end-to-end, attention-free answer path combining directed
  multiworld lenses, finite-speed admissibility, typed GeoTransport, geometric
  memory, min-plus geodesic closure, terminal boundary measurement, and
  fail-closed output;
  \item a controlled multi-seed protocol that separates learnability from
  long-horizon extrapolation while comparing \ERGT with a strengthened direct
  Transformer at the same raw-input boundary; and
  \item frozen causal interventions and a self-contained reproduction package
  that make the roles of geometry, transport, memory, boundary conditions, and
  multiworld composition independently auditable and reproducible.
\end{enumerate}

The claim remains deliberately bounded. \ERGT is not offered as a universal
replacement for Transformers, a general language model, or a cheaper sequence
architecture. The study tests whether geometry induced from raw tokens can act
as an executable and mechanistically testable bias for controlled long-horizon
relational reasoning. Section~\ref{sec:results} reports the corresponding
outcomes, uncertainty estimates, and effect sizes.

\section{Positioning and Related Work}
\label{sec:positioning}

\subsection{Attention and attention-free sequence backbones}

Removing attention is not, by itself, a new research direction. Self-attention
computes content-conditioned interactions across sequence positions
\cite{vaswani2017attention}, whereas S4 realizes long-range propagation through
a structured state-space system \cite{gu2022s4}. Mamba makes that state-space
dynamics input-selective \cite{gu2023mamba}, and Hyena uses implicit long
convolutions with data-controlled gates \cite{poli2023hyena}. These models show
that a sequence can be processed effectively without QKV attention.

Their main object is still the sequence operator: recurrence, state-space
dynamics, or convolution determines how information is mixed over position.
\ERGT addresses a different object. It first constructs a token-indexed field
and asks that field to support a directed, executable relation geometry. Only
then can typed state propagate toward a boundary decision. The contribution is
therefore not attention removal alone, nor a claim about language-modeling
quality or linear-time efficiency, but a raw-input geometry whose constraints
participate directly in inference.

\subsection{Graph reasoning, path methods, and latent structure}

Graph networks expose relational bias through message passing
\cite{battaglia2018relational,gilmer2017neural}, and geometric deep learning
provides the broader language for computation on non-Euclidean structure
\cite{bronstein2021geometric}. Usually, however, the relevant nodes and edges
arrive with the input. NBFNet is especially close in spirit because it
represents a node pair through paths and learns a generalized Bellman--Ford
computation \cite{zhu2021nbfnet}; its paths still belong to a graph that is
already given.

Other work learns structure itself. Neural Relational Inference, for example,
infers an interaction graph while predicting dynamics from observations
\cite{kipf2018nri}. The claim for \ERGT is therefore more specific than
``learning a graph.'' Structure must arise from raw tokens and immediately
become executable through directed metrics, path action, typed transport,
admissibility, memory, and terminal measurement. The same distinction defines
the role of $\Psi$: it is not an extra embedding layer but the audited source
from which identity, pairwise evidence, and every world geometry must descend.
A supplied edge, external entity identifier, or target graph would break that
contract.

\subsection{Neural algorithmic reasoning and size generalization}

The CLRS benchmark made size generalization a central test of learned algorithm
execution \cite{velickovic2022clrs}. Hint-ReLIC likewise shows why intermediate
computation matters, regularizing states to remain invariant under
interventions that should preserve the next algorithmic step
\cite{bevilacqua2023causal}. Both lines of work caution against treating final
accuracy as sufficient evidence of a reusable computation.

\ERGT adopts that evaluation principle without imitating a supplied trace or
set of hints. It learns its geometric dynamics from the sequence and tests
whether the same local transport and admissibility rules continue to compose at
greater depth. After checkpoint selection, named geometric quantities are
removed, shuffled, or restricted while the parameters remain frozen. The
resulting change asks whether a quantity was needed by the learned answer path,
not merely whether supervising it helped during training.

\subsection{Physics-informed inductive bias and information geometry}

Physics-informed models illustrate the value of turning a desired law into an
architectural restriction rather than hoping it appears as a correlation.
Hamiltonian Neural Networks encode conservation in learned dynamics
\cite{greydanus2019hamiltonian}, while explicit Hamiltonian and Lagrangian
constraints extend the same idea \cite{finzi2020constraints}. Geometric deep
learning uses symmetry and geometry to restrict the hypothesis space
\cite{bronstein2021geometric}; information geometry supplies the language of
metrics and dual structure on statistical manifolds \cite{amari2016information}.
\ERGT adopts this as a geometric design analogy rather than a direct
statistical-manifold construction.

\ERGT borrows this design logic, not a literal physical interpretation. Its
terms have concrete computational meanings: $\Psi$ is persistent token
information, worlds are directed relation lenses, action is traversal cost,
the cone is a propagation constraint, and GeoTransport moves typed state. A
geodesic is the least-cost admissible route, conservation protects payload
identity, and boundary measurement determines acceptance. Classical
shortest-path recursion and dynamic programming motivate the closure
\cite{dijkstra1959note,bellman1957dynamic}.

Spectral graph quantities \cite{chung1997spectral} and a Forman-style curvature
proxy \cite{forman2003bochner} are kept outside that answer path. They describe
how the induced geometry changes with horizon and motivate future work on
curvature-guided routing and regularization \cite{topping2022understanding}, as
well as adaptive worlds; they are not used to attribute the current
performance.

\subsection{Position of the claimed contribution}

Table~\ref{tab:positioning} brings these boundaries together. The decisive
difference is not any single familiar component, but the object the model must
construct and execute.

\begin{table}[t]
\centering
\caption{Position of \ERGT relative to neighboring research families.}
\label{tab:positioning}
\small
\begin{tabularx}{\linewidth}{@{}P{0.18\linewidth}P{0.34\linewidth}Y@{}}
\toprule
Family & Primary computational object & Boundary with \ERGT \\
\midrule
Transformer
  & Content-conditioned token mixing
  & Does not explicitly require admissible paths, typed state conservation, or
    fail-closed boundary evidence. \\
S4, Mamba, Hyena
  & State-space recurrence or long convolution over a sequence
  & Replace the sequence operator; they do not by themselves construct and
    execute a relation geometry from raw tokens. \\
GNNs and NBFNet
  & Message passing or path aggregation on a graph
  & Usually receive the graph; \ERGT must induce support and metric before
    transporting typed state. \\
Latent graph inference
  & Inferred interaction structure and predictive dynamics
  & \ERGT additionally turns induced structure into multiworld action,
    admissibility, transport, memory, and a boundary decision. \\
CLRS and Hint-ReLIC
  & Learned algorithm execution and invariant intermediate computation
  & \ERGT learns and intervenes on a native geometric dynamics rather than a
    supplied algorithm trace. \\
Physics-informed learning
  & Dynamics constrained by physical laws or symmetry
  & \ERGT translates physics-inspired principles into operational constraints
    for relational computation, without asserting a literal physical system. \\
\ERGT
  & A raw-token $\Psi$ field and executable induced geometry
  & Unifies structure induction, typed transport, geodesic composition, and
    fail-closed measurement in one raw-input-to-answer path. \\
\bottomrule
\end{tabularx}
\end{table}

No isolated ingredient carries the novelty claim: shortest paths, graph
learning, conservation, attention-free propagation, and boundary conditions
all have established precedents. What is tested here is their integration under
one auditable contract. Raw tokens form a persistent $\Psi$ field; the field
induces several directed geometries; hard constraints govern typed transport;
an admissible geodesic is composed; and only terminal measurement can accept an
answer. This geometry-first contract, rather than general superiority on
sequence tasks, defines the proposed advantage and the empirical burden of the
paper.

\section{Method: From Raw Tokens to Executable Geometry}

\subsection{Conceptual design}

The geometry-first idea becomes useful only if it defines the entire route to
an answer. That route has three responsibilities. It must infer relational
structure from raw symbols without an external graph or identity map; compose
local relations without losing the entity, operator, or payload in transit;
and distinguish a supported terminal state from a candidate that is merely
reachable.

These responsibilities lead to three connected operations. \emph{Geometry
formation} builds the persistent $\Psi$ field, measures pairwise evidence, and
induces several directed world metrics. \emph{Geometric execution} converts
those metrics into admissible actions and moves typed state through their
product space. \emph{Boundary adjudication} accepts the resulting path only if
the correct register, conserved payload, and terminal conditions arrive
together. Geometry is therefore part of the decision process, not an auxiliary
visualization: the output of each operation is the state or constraint required
by the next.

Table~\ref{tab:designlogic} relates the needs of a long relational chain to the
corresponding architectural objects. The physics-derived vocabulary is used in
an operational sense. $\Psi$ denotes persistent token information, a world is
a directed relation lens, action is path cost, the cone masks invalid
propagation, GeoTransport transfers typed state, a geodesic is the least-cost
valid path, and boundary measurement is the final acceptance test.

\begin{table}[t]
\centering
\caption{Conceptual derivation of the architecture.}
\label{tab:designlogic}
\small
\begin{tabularx}{\linewidth}{@{}YYY@{}}
\toprule
Problem need & Design principle & Architecture component \\
\midrule
Preserve token identity and content & Persistent state from raw input &
$\Psi$ information field \\
Represent several relation views & Complementary directed metrics &
Eight world lenses \\
Block impossible propagation & Explicit admissibility &
Finite-speed cone \\
Keep operators bound to state & Type-preserving transfer &
GeoTransport and payload registers \\
Compose long paths & Dynamic programming over action &
Min-plus geodesic closure \\
Reject unsupported candidates & Hard terminal constraints &
Boundary measurement and abstention \\
\bottomrule
\end{tabularx}
\end{table}

These objects are constitutive design choices inspired by physics, not claims
of literal physical law. Their role is testable: equations define them,
invariants monitor their numerical behavior, and frozen interventions reveal
whether the trained model needs them. Continuous neural states construct the
field and local evidence; admissibility, typed-register, and terminal contracts
then decide whether those local pieces form a valid answer. A high answer score
cannot override an invalid intermediate program.

The target inference path is
\[
\begin{aligned}
\text{raw symbols}&\rightarrow\Psi\rightarrow\text{world metrics}
\rightarrow\text{typed event support}\\
&\rightarrow\text{geodesic closure}\rightarrow\text{typed transport}
\rightarrow\text{terminal measurement}.
\end{aligned}
\]
No standard QKV self-attention, Transformer encoder, discrete target compiler,
generic graph executor, or learned answer MLP appears in this path.

\subsection{Raw input and the persistent \texorpdfstring{$\Psi$}{Psi} field}

The construction begins with a state rich enough to support relations without
smuggling in the relation graph. Semantic content alone is not sufficient:
repeated or similar symbols must remain distinguishable, yet identity cannot be
an external pointer to the answer. The initial field therefore joins semantic
content and position with an identity fibre computed from the same raw token.
Every later geometric quantity descends from this shared initial condition.

For token $i$, raw symbol fingerprints are mapped to a train-only semantic
state $x_i$ and a collision-audited identity fibre $u_i$. Both are derived
inside \ERGT from the public raw-token tensors. The initial condition is
\begin{equation}
\Psi_i^{(0)} = f_\psi\!\left(x_i+p_i+\gamma f_{\mathrm{id}}(u_i)\right),
\label{eq:psi}
\end{equation}
where $p_i$ is the position state. The identity term changes both the initial
field and the metrics derived from it, but supplies no pointer, graph edge,
condition label, or answer channel. At the public boundary, the model sees only
raw token identifiers and a padding mask. Equation~\ref{eq:psi} captures the
field-construction claim: relational structure must emerge from pairwise
observables of the same persistent state built from the serialization seen by
the baseline.

\subsection{Field observables and adhesion}

The token field still has to determine which pairs can support a relation.
Rather than compressing that decision into one unconstrained edge logit,
\ERGT exposes several bounded factors: coherence, reconstructibility, local
compatibility, order, and agreement with the bulk--boundary view. Together they
turn the node-like state $\Psi$ into candidate edge support while leaving the
source of weak support visible.

The hidden field induces the bounded observables
\begin{align}
\rho_i &= d^{-1}\|\Psi_i\|_2^2, \\
C_{ij} &= \tfrac12\!\left[1+\langle f_C(\Psi_i),g_C(\Psi_j)\rangle\right], \\
R_{ij} &= \sigma\!\left(2\langle f_R(\Psi_i),g_R(\Psi_j)\rangle\right), \\
G_{ij} &= \exp(-|\rho_i-\rho_j|), \\
B_{ij} &= \exp[-(\delta_i^{\mathrm{bulk}}+
                         \delta_j^{\mathrm{bulk}})/2].
\end{align}
$C$ measures coherence, $R$ reconstructibility, $G$ density-gradient
compatibility, and $B$ bulk--boundary agreement. Local relational entropy
$S_i$ gives the order term $O_i=1-S_i$, while stability is measured across
field updates. The factors meet in an adhesion field formed by a geometric mean
in log space:
\begin{equation}
\Phi_{ij}=\exp\!\left(\frac{1}{6}\sum_{k=1}^{6}
                         \log(z_{ij}^{(k)}+\epsilon)\right),
\quad
z^{(k)}\in\{C,R,\mathrm{stability},G,O,B\}.
\label{eq:adhesion}
\end{equation}
Because the factors combine multiplicatively, one weak term cannot be hidden
by several strong ones. $\Phi$ is an edge-formation gate chosen for this model,
not a universal law of information.

\subsection{Constitutive world geometries}

Adhesion establishes joint support, but not a single universal notion of
proximity. Locality, reconstruction, long bridges, density flow, and boundary
compatibility may favor different directed edges. Collapsing them immediately
would erase those alternatives, so \ERGT applies eight constitutive lenses to
the shared field and keeps their potentials separate until a typed role needs
them. Table~\ref{tab:worlds} summarizes the resulting worlds.

\begin{table}[t]
\centering
\caption{Operational world lenses. These are constitutive functions, not
independently supplied graphs.}
\label{tab:worlds}
\small
\begin{tabularx}{\linewidth}{@{}lY@{}}
\toprule
World & Dominant factors \\
\midrule
$W_0$ & locality, density-gradient fit, stability \\
$W_1$ & coherence, identity overlap, local order \\
$W_2$ & reconstruction and boundary reconstructibility \\
$W_3$ & stability, identity overlap, locality \\
$W_4$ & long bridge, coherence/identity support, reconstruction \\
$W_5$ & density contrast and stability \\
$W_6$ & novelty and long-bridge support \\
$W_7$ & density source--sink flow and coherence \\
\bottomrule
\end{tabularx}
\end{table}

For world $w$, edge state evolves as
\begin{equation}
E_{wij}^{(t+1)}=\beta E_{wij}^{(t)}+
 (1-\beta)\Phi_{ij}^{(t)}\Omega_{wij}^{(t)},
\qquad
\ell_{wij}=-\log(E_{wij}+\epsilon).
\label{eq:edge}
\end{equation}
Typed computational roles condition a router over these separate metrics. Let
$\bar\Psi$ be the masked mean field, $r$ a typed role, $\pi_{rw}$ a learned
role--world prior, and $\mathcal{W}_r$ the worlds authorized for that role:
\begin{align}
\alpha_{rw} &= \operatorname{softmax}_{w\in\mathcal{W}_r}
   \left([W_{\mathrm{route}}\bar\Psi]_{rw}+\pi_{rw}\right), \\
H_r &= \sum_{w\in\mathcal{W}_r}\alpha_{rw}A_w(H_w),
\label{eq:router}
\end{align}
where $A_w$ adapts world $w$. A world is therefore not a separate predictor,
but one directed metric through which the shared typed state can be viewed and
transported. The worlds need not carry equal mass or have unique
human-readable meanings. The testable claim is functional dependence on their
composition: under frozen intervention, no single world completes the program.

\subsection{Finite-speed admissibility and action}

An edge potential says where a relation appears plausible; it does not yet make
that edge executable. Keeping every nonzero affinity would leave weak, dense
shortcuts through which the model could bypass a long relation chain. We
therefore turn each learned world geometry into a sparse feasible domain and
assign a non-negative cost to every legal local move. Admissibility determines
\emph{whether} state may propagate, while action distinguishes among the routes
that remain.

Sparse support and a finite propagation condition define admissible edges:
\begin{equation}
\mathcal{A}_{wij}=
\mathbf{1}[j\in\operatorname{topk}(E_{wi\cdot})]
\mathbf{1}[\ell_{wij}\le v_i\Delta\tau_{ij}].
\end{equation}
The path action is
\begin{equation}
S_{wij}=\ell_{wij}+\lambda_b(1-B_{ij})+
\lambda_s(1-\mathrm{stability}_{ij})+
\lambda_c\,\mathrm{coneViolation}_{ij}.
\label{eq:action}
\end{equation}
This action is a learned non-negative path cost. The term \emph{finite-speed
cone} means that the implemented update rejects edges whose effective travel
cost exceeds the available discrete propagation interval; it is not a claim
of relativistic spacetime.

\subsection{GeoTransport and memory}

A path cost is useful only if information can travel along the path and remain
available to later steps. GeoTransport provides that state evolution. Its mask
is inherited from admissibility and its weights from action, so an edge rejected
by the geometric stage cannot reappear during transfer. Persistent momentum,
edge state, and payload retain the history on which later geometry and terminal
measurement depend.

Transport is induced by action over admissible neighbors:
\begin{align}
T_{wij} &= \operatorname{softmax}_{j:\mathcal{A}_{wij}=1}
                  (-S_{wij}/\tau), \\
M_{wi} &= \sum_j T_{wij}V_w(\Psi_j), \\
\Pi_i^{(t+1)} &= \mu\Pi_i^{(t)}+(1-\mu)
                  F(\Psi_i^{(t)},M_i^{(t)}), \\
\Psi_i^{(t+1)} &= \operatorname{Norm}
                  (\Psi_i^{(t)}+\Pi_i^{(t+1)}).
\end{align}
Memory is path-dependent state retained in $E$, $\Pi$, and transported payload.
It can affect later geometry and terminal measurement but has no direct bypass
to answer logits. In computational language, GeoTransport is therefore a
masked, action-weighted transfer operator that keeps a typed state attached to
the path on which it arrived; it is not attention-weighted value mixing.

\subsection{Typed event support}

Geometry determines where state may travel and what the journey costs, but it
does not by itself identify the relation operator carried by an event. That
semantic content must remain sufficiently discrete to prevent a short route
from substituting one operator for another. Typed event support therefore binds
source, relation, and target proposals to the path, while the geometry retains
control over admissibility and route selection. This provides the interface
between learned interpretation of the raw symbols and native geometric
execution, without introducing a compiled target graph.

Predicted event source, relation, and target distributions form typed charge
\begin{equation}
Q_{ijr}=\sum_e p(i\mid e)\,p(r\mid\Psi_e^{(0)})\,p(j\mid e).
\end{equation}
The charge supplies typed local proposals to the geometric program. Geometry
determines whether and at what cost a proposal can participate in a complete
path; it cannot transmute the discrete operator type carried by $Q$. This
separation prevents a low scalar path cost from silently changing the relation
being transported.

\subsection{Product-manifold geodesic closure}

The preceding stages define legal one-step transitions. To reason over a long
horizon, the model must compose them without handing the result to a fixed-depth
answer classifier. \ERGT does so by running dynamic programming over the
discrete product of token location and world lens. One coordinate records where
the state is; the other records the directed metric under which its current
transition is evaluated. Min-plus closure can then extend the same local action
contract to paths that exceed a single field update.

The solver state is the discrete product $(i,w)$ of token and world. For
candidate $c$, the min-plus field is initialized at its registered source:
\begin{align}
D_c^{(0)}(i,w) &=
\begin{cases}0,&(i,w)\text{ is an allowed source},\\
+\infty,&\text{otherwise},\end{cases}\\
D_c^{(t+1)}(j,w') &= \min\left\{D_c^{(t)}(j,w'),
D_c^{(t)}(i,w)+A[(i,w)\rightarrow(j,w')]\right\}.
\label{eq:minplus}
\end{align}
Support is propagated by a max--min widest-path channel so repeated
probability multiplication alone cannot erase a valid long path. Closure depth
is at least the registered answer horizon and is capped at 36 steps in the
paper profile. A label-free dyadic valid-domain mesh allows action to
propagate over long raw sequences without supplying the answer topology.
Here \emph{geodesic} has a precise discrete meaning: the path of minimum total
implemented action among paths that satisfy the admissibility mask.

\subsection{Typed payload and terminal decision}

A least-action route establishes reachability, but reachability alone is not an
answer. The arriving state must also encode the operator sequence required by
the candidate, preserve enough transported support, and satisfy the registered
terminal boundary. The final stage couples path optimization to typed register
execution and asks, through hard measurements, whether the assembled evidence
is complete. This is where fail-closed behavior becomes part of the computation
rather than an output convention.

Along an admissible event, the register and payload evolve as
\begin{align}
r' &= \operatorname{applyOperator}(r,r_e),\\
q^{(t+1)}(j,r') &= \max_e q^{(t)}(i_e,r)\,\tau_e m_e,\\
\Delta^{(t+1)}(j,r') &= \min_e
             \left[\Delta^{(t)}(i_e,r)+\delta_e\right].
\end{align}
Two normalization contracts are checked before checkpoint eligibility:
\begin{align}
\varepsilon_{\mathrm{typed}}
  &=\max_e\left|\sum_r p_e(r)-p_e(\mathrm{present})\right|,\\
\varepsilon_{\mathrm{transport}}
  &=\max_{w,i\ \mathrm{valid}}\left|\sum_j T_{wij}-1\right|.
\label{eq:conservation}
\end{align}
The first keeps relation-type mass tied to event-presence mass; the second
keeps each valid transport row normalized. For candidate $c$, let $A_c$ be its
closed geodesic action, $m_c$ its terminal payload mass, and $\delta_c$ its
boundary deficit. Its hard validity indicator is
\begin{equation}
V_c=\mathbf{1}[A_c<\infty]
    \mathbf{1}[m_c\ge m_{\min}]
    \mathbf{1}[\delta_c\le\delta_{\max}],
\label{eq:validity}
\end{equation}
where $A_c$ is read from the register state required by that candidate, so a
wrong typed register is unsupported. Cone validity is already enforced in the
edges from which finite $A_c$ can be assembled.
The decision is
\begin{equation}
\widehat c=\arg\min_{c:V_c=1}A_c
\quad\text{only if}\quad
A_{(2)}-A_{(1)}\ge\eta;
\quad\text{otherwise abstain}.
\label{eq:decision}
\end{equation}
Thus boundary measurement means checking the final register, mass,
conservation, and deficit before accepting a candidate. Two unresolved
candidates, no valid candidate, or insufficient action margin produces
abstention. This is a native geometric boundary solver, not a call to a generic
graph-traversal routine.

\subsection{Training objective and zero-teacher verification}

Training must reconcile differentiable field construction with hard execution
contracts. The upstream field needs enough supervision to populate the native
path, yet the reported inference cannot depend on a teacher or a soft answer
bypass. We therefore group intermediate losses by architectural role and keep
checkpoint selection separate from the subsequent frozen execution test.

The grouped training objective is
\begin{equation}
\mathcal{L}=\mathcal{L}_{\mathrm{answer}}
+2a(t)\mathcal{L}_{\mathrm{structured}}
+2a(t)\mathcal{L}_{\mathrm{physical}}
+2a(t)\mathcal{L}_{\mathrm{causal}}
+2a(t)\mathcal{L}_{\mathrm{margin}}
+0.5\mathcal{L}_{\mathrm{geometry}}.
\label{eq:loss}
\end{equation}
$a(t)$ is a teacher schedule. Structured terms supervise entity, query role,
event type, endpoints, operator, boundary, and edge assembly. Physical terms
supervise action, cone, transport, boundary deficit, and terminal mass.
Geometry terms include bulk--boundary reconstruction, conservation, world
usage, and diversity. A checkpoint becomes eligible only after independent
field-readiness windows. Teacher weight is then set to zero, parameters are
frozen, and two further windows must pass with zero optimizer steps. Thus the
reported answer is teacher-free at inference, although training is supervised.

\section{Experimental Design}
\label{sec:experimental}

The experiment mirrors the computational chain defined in Section~3. It first
asks whether raw symbols can produce a stable native geometry, then whether the
frozen computation survives longer relation chains, and finally whether its
answers depend on the geometric operations claimed by the method. Unsupported
inputs add a separate test of whether the terminal rule can withhold an answer.
These questions cannot be answered by a single accuracy score, so each is
assigned its own evidence panel.

The panels are also kept separate during model selection. Development data are
used only to lock the strengthened direct-Transformer recipe; native readiness
data select the first stable \ERGT checkpoint; and the final horizon cohorts
remain unopened until both procedures are complete. Interventions are applied
only to that frozen native checkpoint. Table~\ref{tab:studylogic} states the
resulting inferential boundary.

\begin{table}[t]
\centering
\caption{Inferential structure of the registered study.}
\label{tab:studylogic}
\small
\begin{tabularx}{\linewidth}{@{}P{0.25\linewidth}P{0.34\linewidth}Y@{}}
\toprule
Question & Evidence panel & Permitted conclusion \\
\midrule
Was the inference path learned?
  & Train ID, raw-disjoint validation, and native stability contracts
  & Model-specific convergence and, for \ERGT, executable-path readiness \\
Does computation extend in horizon?
  & Frozen matched cohorts beyond the training range
  & Bounded long-horizon accuracy and paired model difference \\
Does the model fail closed?
  & Controlled inputs with no valid candidate
  & Rejection within the registered unsupported protocol \\
Is geometry functionally necessary?
  & Same-checkpoint removals, corruption, and single-world controls
  & Causal attribution to the intervened mechanism on eligible solved pairs \\
How does geometry change with horizon?
  & Detached curvature, spectrum, and world-diversity observers
  & Descriptive structure only, not causal performance attribution \\
\bottomrule
\end{tabularx}
\end{table}

\subsection{Benchmark and counterfactual cohorts}

The benchmark isolates relational composition from open-domain language
knowledge. Each example serializes nodes, relation events, a query, two
candidates, and distractors as raw symbols; this serialization is the complete
input to both models. A latent graph exists only inside the generator so that
examples can be constructed and matched. Adjacency, edge indices, entity
pointers, and candidate paths are never exposed.

The basic unit is an answer-reversing counterfactual pair. Its members share
topology, edge count, hop count, operator types, and lexical statistics, but one
registered condition changes the valid candidate. Scenarios separately vary
geodesic action, finite-speed admissibility, boundary deficit, payload
transport, terminal mass, source memory, and their combined long-horizon use.
The pair design therefore rewards tracking the condition that governs a path,
not memorizing its surface form or candidate position.

Training contains 28 pair groups at each hop from 1 through 8: 224 supported
groups in total, plus 28 controlled unsupported groups. Independent native
readiness data cover hops 2--8. Frozen evaluation uses 4, 8, 12, 16, 24, and
32 hops; context panels use 128, 256, 512, and 768 tokens; and unsupported
cases are tested at 4, 16, and 32 hops. The horizon panels vary relational
depth, whereas the context panels add serialized material around a fixed
decision. This controlled synthetic setting is intended to isolate the
mechanism and does not measure open-domain language ability.

\subsection{Compared inference paths}

Fairness is defined at the external task boundary rather than by forcing the
architectures to share an internal solver. Both receive the same serialization
and padding mask, train on the registered 1--8-hop family under locked
schedules, and predict candidate A, candidate B, or abstention. The baseline is
strengthened during training, but no ERGT-derived structure enters its inference
path. The comparison is therefore between two complete ways of deriving an
answer from raw symbols, not between \ERGT and a Transformer equipped with the
\ERGT solver.

The direct arm is a standard PyTorch Transformer encoder with token embeddings,
sinusoidal positions, QKV multi-head attention, GELU feed-forward blocks,
layer normalization, and a three-class answer head. Its locked configuration
uses width 64, four heads, four layers, feed-forward width 192, context 768,
and zero dropout. It has 205,634 inference-active parameters, compared with
196,816 for \ERGT, a 4.48\% difference. Matched auxiliary targets strengthen
training but are not inference inputs, and the model receives no graph,
geometric state, condition side input, compiler, program, executor, or identity
side channel.

Native \ERGT uses width 56, $\Psi$ rank 20, eight world lenses, three field
updates, sparse top-$20$ support, and at most 36 geodesic-closure steps. From
the same serialization and mask, it must construct identity, relation support,
world geometry, typed transport, and the terminal decision through the path in
Section~3; no learned answer MLP bypasses that route. Parameter proximity limits
a simple capacity-scale explanation, but it does not imply equal computation.
Optimizer steps, time, and peak CUDA allocation are therefore measured rather
than matched.

\subsection{Locked training and checkpoint selection}

The Transformer recipe was qualified once on development seeds before any
12--32-hop final cohort was opened. It begins with one-hop warm-up at batch
size 8 and learning rate $7\times10^{-4}$; after two eligible raw-disjoint ID
windows, it may enter fixed stages with exact 1--2, 1--4, and 1--8-hop
coverage. Architecture, optimizer, transition rule, and step budget are then
locked. A missed transition is recorded rather than repaired using final data.

The confirmatory \texttt{four\_seed\_confirmation} profile trains both models
from scratch. Optimization seeds 17117, 18313, 19501, and 20707 are paired
with independent data seeds 21101, 22307, 23509, and 24733. All raw sets are
collision-audited and disjoint across training, tuning, validation, and final
evaluation. Horizon, context, unsupported, intervention, and observer panels
remain outside checkpoint selection.

During confirmation, the Transformer follows its locked warm-up and curriculum
within 1400 steps. \ERGT uses the registered 1--8-hop pool, batch size 8,
AdamW at $2.5\times10^{-3}$, and at most 1800 steps. Its teacher is used only
during training. The first checkpoint to complete two 2--8-hop field-readiness
windows is frozen immediately, then evaluated in two unchanged zero-teacher
windows without another optimizer update. Selection therefore depends on a
stable short-to-intermediate native program, never on the final horizon score.

\subsection{Registered outcomes and interventions}

Row accuracy records individual answers; counterfactual-pair exactness requires
both members of an answer-reversing pair to be correct. Model-specific
convergence requires at least 0.90 supported accuracy on both train ID and
untouched raw-disjoint validation. A secondary matched-convergence analysis is
restricted to seed pairs for which both models satisfy that rule. Failure in one
arm limits only joint-convergence claims and does not erase the other model's
independent result.

The preregistered primary comparison is the 32-hop panel. Its bounded claim
requires native accuracy of at least 0.90, a margin of at least 0.35 over the
direct Transformer, a hierarchical bootstrap interval excluding zero, and
exact paired McNemar $p<0.05$. Counterfactual pairs are the paired clusters and
are nested within data seed. Bootstrap intervals follow Efron and Tibshirani
\cite{efron1993bootstrap}; discordant outcomes use McNemar's exact test
\cite{mcnemar1947note}. Intermediate horizons, context lengths, pair exactness,
and unsupported rejection remain supporting panels and cannot select a
checkpoint.

Native checkpoint eligibility is stricter than answer accuracy. Measurement
accuracy must reach 0.99; action MAE must be at most 0.12; event recall,
precision, path coverage, and chain exactness must reach 0.99, 0.98, 0.99,
and 0.99; closure must not degrade; hard overflow must be zero; conservation
residual must be at most $10^{-4}$; every group must have exposure 8; and two
consecutive readiness windows must pass. Batch, padding, identity, and raw-input
invariance are recorded with these contracts. All registered outcomes remain
in the readout even when a claim-specific rule is not met.

Mechanism attribution then reuses the selected checkpoint without retraining.
It changes one inference operation at a time and evaluates only
counterfactual pairs solved by the intact model. Targeted removals disable
action, cone, typed transport, boundary deficit, terminal mass, or
source-conditioned memory; global controls randomize or shuffle geometry,
bypass the conditioned world gate, or retain one world. The minimum targeted
drop is 0.20. Formal evaluation requires intact accuracy of at least 0.80 and
at least two eligible pairs; conditional effects are still reported below that
panel-wide floor, but the corresponding claim is narrowed.

Curvature, graph gradient, Laplacian spectrum, spectral entropy, effective
rank, spectral gap, and world diversity are handled differently. They observe
the frozen geometry outside the answer path and cannot change a prediction.
Their role is descriptive and prospective, not causal attribution for current
performance.

\section{Results}
\label{sec:results}

The evidence follows the life of a single learned computation. We first verify
that the native path converges and can be frozen without losing its contracts.
That fixed path is then carried to greater relational depth and longer contexts
before the same checkpoints face unsupported inputs, mechanism interventions,
and detached geometric observation. Fitting, extrapolation, and causal
attribution thus remain separate empirical questions, but their answers refer
to the same computation.

\subsection{Convergence and checkpoint integrity}

The long-horizon study begins with the computation that will be extrapolated.
As Table~\ref{tab:convergence} shows, native \ERGT reached 1.000 accuracy on
both training and raw-disjoint validation in all four optimization/data-seed
pairs. The complete route---from $\Psi$ formation through world geometry and
GeoTransport to terminal measurement---therefore converged across fresh data
and initializations. The strengthened direct Transformer also reached 1.000 on
every training split, but its held-out accuracy ranged from 0.449 to 0.837.
This gap describes the registered baseline; it does not alter the independent
convergence result for \ERGT.

\begin{table}[t]
\centering
\caption{Accuracy by independent optimization/data-seed pair. Values are
from fresh confirmatory runs.}
\label{tab:convergence}
\small
\begin{tabular}{@{}lccccc@{}}
\toprule
Seed pair & \ERGT train & \ERGT val. & \TF train & \TF val. & Freeze step \\
\midrule
17117/21101 & 1.000 & 1.000 & 1.000 & 0.510 & 700 \\
18313/22307 & 1.000 & 1.000 & 1.000 & 0.459 & 700 \\
19501/23509 & 1.000 & 1.000 & 1.000 & 0.837 & 1000 \\
20707/24733 & 1.000 & 1.000 & 1.000 & 0.449 & 1000 \\
\bottomrule
\end{tabular}
\end{table}

Accuracy was not sufficient to select a checkpoint. The first state to pass two
complete readiness windows was frozen only after every registered 2--8-hop
group had sufficient exposure. The four freeze steps were 700, 700, 1000, and
1000; minimum group exposures were 11, 11, 15, and 15 against the required
floor of 8. No group was unseen, and no optimizer update followed freeze.

The internal audits reach the same conclusion at finer resolution. All 56
instances of the 14 native stability invariants passed. Event-chain exactness,
geodesic path coverage, closure non-degradation, event precision, and event
recall were each 1.000; overflow was zero; and the largest typed-payload
conservation residual was $2.38\times10^{-7}$. Batch-size and padding
invariance, identity construction from raw input, and raw-input parity also
passed. The selected states were therefore stable executions of the intended
native path rather than products of checkpoint drift, an external identity
channel, or serialization artifacts.

\subsection{Frozen generalization across horizon and context}

Those states were then held fixed throughout the horizon study. Because
training groups covered 1--8 hops and no final panel could influence selection,
Table~\ref{tab:horizon} measures extrapolation rather than adaptation. Mean
\ERGT accuracy remained above 0.90 at every registered horizon and reached
0.9554 at 32 hops, whereas the direct-Transformer means remained near 0.50.

\begin{table}[t]
\centering
\caption{Mean accuracy across four independent optimization/data-seed pairs.
Training groups cover 1--8 hops.}
\label{tab:horizon}
\begin{tabular}{@{}rccr@{}}
\toprule
Hops & Native ERGT & Direct Transformer & Difference \\
\midrule
4  & 1.0000 & 0.4821 & 0.5179 \\
8  & 1.0000 & 0.5045 & 0.4955 \\
12 & 0.9821 & 0.5000 & 0.4821 \\
16 & 0.9866 & 0.5000 & 0.4866 \\
24 & 0.9196 & 0.5000 & 0.4196 \\
32 & 0.9554 & 0.5000 & 0.4554 \\
\bottomrule
\end{tabular}
\end{table}

The aggregate result includes visible seed variation. At 32 hops, \ERGT
accuracy was 0.982, 1.000, 1.000, and 0.839 across the four runs; exactness on
the corresponding answer-reversing counterfactual pairs was 0.964, 1.000,
1.000, and 0.750. The pooled accuracies are
\[
\operatorname{Acc}_{\ERGT}=0.955357,
\qquad
\operatorname{Acc}_{\TF}=0.500000.
\]
Their paired difference, 0.455357, exceeds the registered 0.35 margin, while
the native mean exceeds its 0.90 floor. A hierarchical bootstrap over data
seed and counterfactual pair gives the 95\% interval $[0.375,0.500]$, and the
exact paired test gives
\[
p_{\mathrm{McNemar}}=9.86\times10^{-32}.
\]
Of the 224 examples in 112 pairs, 104 were correct only for \ERGT and none
only for the Transformer, matching the direction of the bootstrap interval and
exact test. Within this benchmark, the frozen geometric computation therefore
retains substantially more accuracy as relational depth extends to 32 hops.

The context panels distinguish that effect from tolerance of extra tokens. At
hop 4, mean \ERGT accuracy was 1.000, 1.000, 1.000, and 0.963 at 128, 256,
512, and 768 tokens, compared with 0.500, 0.506, 0.519, and 0.500 for the
Transformer. At hop 24 and contexts of 512 and 768 tokens, the corresponding
values were 0.931 and 0.988 for \ERGT versus 0.519 and 0.531. The horizon
effect thus persists inside the registered longer serializations, without being
presented as a general long-context language-modeling result.

\subsection{Fail-closed behavior and geometric attribution}

The next question is not how often the model answers, but whether it answers
only when its own terminal conditions are satisfied. On controlled cases in
which neither candidate has admissible transported support, \ERGT rejected
every example at 4, 16, and 32 hops in every seed
(Table~\ref{tab:unsupported}); the direct Transformer rejected none.

\begin{table}[t]
\centering
\caption{Mean correct rejection rate on controlled unsupported candidates.}
\label{tab:unsupported}
\begin{tabular}{@{}rcc@{}}
\toprule
Hops & Native ERGT & Direct Transformer \\
\midrule
4  & 1.0000 & 0.0000 \\
16 & 1.0000 & 0.0000 \\
32 & 1.0000 & 0.0000 \\
\bottomrule
\end{tabular}
\end{table}

This is the benchmark's fail-closed result: absence of valid geometric support
produces abstention rather than a forced choice. It is not an open-domain
hallucination measure.

The same frozen checkpoints were then used to test whether successful answers
actually depend on the proposed geometry. Table~\ref{tab:mechanisms} reports
scenario-matched interventions on pairs solved by the intact model, so each
drop is the loss of an established capability. Action ranks admissible paths,
the cone excludes invalid propagation, transport carries typed state, terminal
mass verifies its arrival, memory preserves source information, and boundary
deficit checks terminal compatibility. Removing action, cone, transport, or
terminal mass reduced eligible accuracy from 1.000 to 0.000 in all four seeds;
removing memory left 0.031.

\begin{table}[t]
\centering
\caption{Frozen mechanism interventions. Full and intervention accuracies are
computed on eligible solved pairs; ``formal seeds'' also satisfy the panel-wide
0.80 intact-model floor.}
\label{tab:mechanisms}
\small
\begin{tabular}{@{}lccc@{}}
\toprule
Removed mechanism & Full & Intervention & Formal seeds \\
\midrule
Geodesic action & 1.000 & 0.000 & 4/4 \\
Finite-speed cone & 1.000 & 0.000 & 4/4 \\
Typed payload transport & 1.000 & 0.000 & 4/4 \\
Terminal mass & 1.000 & 0.000 & 4/4 \\
Source-conditioned memory & 1.000 & 0.031 & 4/4 \\
Boundary deficit & 1.000 & 0.000 & 2/4 \\
\bottomrule
\end{tabular}
\end{table}

Boundary removal likewise changed eligible solved-pair accuracy from 1.000 to
0.000 in every seed. Formal panel-wide qualification was 2/4 because two intact
boundary panels scored 0.750, below the conservative 0.80 evaluability floor.
The conditional effect is therefore present in all four runs, while the formal
panel claim remains limited to the two preregistered eligible seeds.

The global controls lead to the same interpretation. Random geometry reduced
eligible accuracy by 0.646, cross-example shuffling by 0.597, and bypassing the
conditioned world gate by 0.382. Restricting inference to any one of worlds
0--7 produced zero accuracy. The result thus depends on the learned arrangement
of costs and constraints and on composition across worlds, without requiring
equal world use or a unique human-readable meaning for each world.

\subsection{Detached geometric observers}

Curvature and spectrum were observed alongside the answer path, never used as
answer features. The detached CPU/float64 pipeline processed 26,176 induced
graphs without unavailable values, sanitization, numerical failure, or a CUDA
eigensolver call. In Table~\ref{tab:observers}, curvature describes local graph
expansion and contraction; spectral rank, entropy, and gap summarize global
mode use and connectivity; and world diversity records how broadly the
geometric lenses participate.

\begin{table}[t]
\centering
\caption{Detached geometric observers, averaged over four seeds.}
\label{tab:observers}
\small
\begin{tabular}{@{}lcc@{}}
\toprule
Observer & 4 hops & 32 hops \\
\midrule
Mean curvature proxy & $-0.1011$ & $-0.0154$ \\
Negative-curvature fraction & 0.8959 & 0.9479 \\
Mean absolute curvature gradient & 0.002424 & 0.000354 \\
Spectral effective rank & 4.859 & 43.867 \\
Spectral entropy & 0.711 & 0.868 \\
Spectral gap & 0.1512 & 0.0115 \\
World diversity & 0.1447 & 0.8052 \\
\bottomrule
\end{tabular}
\end{table}

Between 4 and 32 hops, effective spectral rank rose from 4.859 to 43.867,
entropy from 0.711 to 0.868, and world diversity from 0.1447 to 0.8052.
Spectral gap fell from 0.1512 to 0.0115 and mean absolute curvature gradient
from 0.002424 to 0.000354, while the Forman-style curvature proxy remained
predominantly negative. The geometry therefore becomes more distributed across
modes and worlds as chains lengthen. Because the observers were detached, this
is a structural description, not evidence that curvature or spectrum caused
the accuracy gain.

\subsection{Computational cost}

The reported behavior carries a higher computational cost. Mean training time
on the Colab runtime was 1437.5 seconds for \ERGT and 56.25 seconds for the
direct Transformer, a per-seed ratio of 21.3--29.8. Mean peak CUDA allocation
was approximately 3.432 GiB versus 1.455 GiB, a ratio of 2.36. These values are
reported for transparency and do not support an efficiency claim.

Read as one sequence, the four-seed evidence connects convergence of the native
raw-token-to-$\Psi$ path to frozen 32-hop generalization, fail-closed behavior,
and mechanism-specific intervention effects. The numerical advantage is thus
paired with evidence that the learned geometry, transport, memory, and terminal
constraints participate in the answer, at the greater resource cost reported
above.

\section{Discussion}
\label{sec:discussion}

The results support a direct conclusion: \ERGT learned the complete route from
raw tokens to a constrained decision. The $\Psi$ field was not an intermediate
representation handed to a conventional readout. It supplied the identity and
relation evidence from which the model formed directed world geometries,
transported typed state, and licensed a terminal answer. This route converged
in every fresh run and remained intact after the checkpoint was frozen. The
experiment therefore establishes more than the feasibility of writing down a
geometric model; it shows that the proposed geometry can be learned end to end
and used as the model's native inference path.

The horizon panels then ask whether that learned computation is reusable. They
were never used to choose a checkpoint, so the performance reported at long
relation depths reflects a frozen rule learned from shorter compositions rather
than adaptation to the test horizon. The context controls help locate the
effect: increasing the number of surrounding tokens is not equivalent to
increasing the number of relations that must be composed. Across the registered
panels, the geometric path remained strongest when the task required deeper
composition, while the unsupported cases showed that it could also withhold an
answer when the terminal conditions were not met.

The direct-Transformer comparison gives this result a clear, limited meaning.
Both models start from the same raw serialization, receive matched training
labels, and have closely matched active parameter counts. The Transformer was
trained with the strengthened, locked recipe selected during development, but
it received no ERGT geometry or solver at inference. It fitted the fresh
training sets yet did not carry the same rule across raw-disjoint validation and
longer horizons. Under this protocol, explicit construction and execution of a
relation geometry generalized better than direct content mixing. This is an
architectural result for the registered benchmark, not a claim that every
Transformer must fail or that the comparison represents universal
post-convergence superiority.

The interventions explain why the accuracy difference is attributable to the
proposed computation. Removing path action, propagation admissibility, typed
transport, terminal mass, or source-conditioned memory broke distinctions that
the intact model solved. Random and shuffled geometries also caused substantial
losses, and no single world reproduced the complete computation. These results
rule out a simple account in which the geometric quantities are merely labels
attached to an unrelated predictor. They support the mechanism as a composed
system: the evidence does not require every component to be sufficient on its
own, but it does show that the registered path depends on its geometry,
transport, memory, and terminal constraints.

This dependence also offers a natural explanation for the long-horizon
behavior. $\Psi$ preserves identity and local relation evidence; the worlds
retain distinct directional views; the cone removes inadmissible propagation;
action ranks the remaining paths; GeoTransport carries the typed payload and
memory; and terminal measurement decides whether the transported state is a
valid answer. Extending a chain repeats these local contracts without changing
their meaning. The observed horizon generalization is therefore consistent
with reuse of a stable geometric law rather than reconstruction of a new global
correlation at each length. The frozen interventions turn that interpretation
into a testable account rather than a metaphor applied after prediction.

The broader contribution is consequently not the absence of attention by
itself. It is the construction of task-relevant structure from raw input and
the immediate use of that structure as an executable, fail-closed geometry.
The model is not given a graph, and the induced graph is not retained only as
an explanation: its support, direction, cost, transported state, and boundary
conditions determine the answer. This places \ERGT between sequence modeling,
latent structure learning, path reasoning, and physics-informed design, while
the combination tested here remains distinct from simply replacing the
sequence mixer or running a learned algorithm on a supplied graph.

Finally, the curvature and spectral measurements describe how the induced
geometry changes with depth, but they do not participate in the present
decision path. Their trends suggest concrete future tests: curvature may help
regularize action or detect transport bottlenecks, while spectral rank, gap,
and entropy may help monitor world allocation, collapse, memory organization,
or structural shift. Those possibilities remain prospective. The result
established here is narrower and firmer: within the registered relational
family, a raw-token information field can support a convergent native geometry,
and that geometry can sustain auditable long-horizon reasoning.

\section{Scope and Limitations}

This study addresses a bounded architectural hypothesis: geometry induced from
raw tokens can execute long-horizon relational reasoning. Internal validity
within the registered synthetic family is supported by four independent seed
pairs, fresh data, frozen checkpoints, paired counterfactuals, invariant audits,
and same-checkpoint interventions. Whether the result extends to natural
language, broader tasks, or other model scales remains open.

The parameter-matched direct Transformer was strengthened and locked before the
confirmatory runs, but did not reproduce its development held-out criterion on
the fresh data. The comparison therefore establishes an architectural advantage
for this benchmark, not universal or post-matched-convergence superiority.
Training uses structured supervision, whereas native inference is teacher-,
compiler-, and executor-free. Boundary dependence is exact on eligible solved
pairs and meets the formal panel-wide criterion in two of four seeds. Curvature
and spectrum remain descriptive observers, computational efficiency is not
claimed, and unsupported rejection refers only to the registered protocol.

\section{Code and Data Availability}

The complete reproduction release is available in the public
\href{https://github.com/jalaljafari2009/ERGT-paper}{\texttt{ERGT-paper}}
repository. To run it, first choose \texttt{Code > Download ZIP} on that page
and extract the downloaded \texttt{ERGT-paper-main.zip}. Then open the
\href{https://colab.research.google.com/github/jalaljafari2009/ERGT-paper/blob/main/notebook/ERGT_Attention_Free_Geometric_Reasoning_Study.ipynb}{executable Colab notebook},
select a GPU runtime, and choose \texttt{Run all}. When the notebook opens the
upload dialog, select the inner archive
\path{ERGT-paper-main/ERGT_FourSeed_Reproduction.zip} from the extracted
folder. The outer \texttt{ERGT-paper-main.zip} is not the runtime package and
must not be uploaded to the notebook.

\section{Conclusion}

\ERGT demonstrates a distinct route from sequence input to relational
reasoning. Instead of using self-attention or a Transformer backbone to mix the
sequence and read out an answer, it organizes raw-token information into a
persistent $\Psi$ field. Identity and relation evidence in that field induce
multiple directed geometries; GeoTransport carries typed state through their
admissible geodesics; and boundary measurement accepts an answer only when the
transported state satisfies the terminal constraints. The induced geometry is
therefore neither a supplied graph nor a post-hoc explanation. It is the
model's native inference machinery.

This complete field-to-answer path converged in four independent runs, passed
all 56 stability invariants, and retained mean accuracy 0.955 at 32 hops,
exceeding the strengthened direct Transformer by 0.455 with strong paired
evidence. Frozen interventions further tied the solved answers to geometric
action, admissibility, transport, memory, terminal mass, and multiworld
composition. Within the registered relational family, these results establish
that a trainable, attention-free information geometry can form from raw input,
converge end to end, and reuse its local transport laws for substantially
longer chains. This is the central contribution and the basis for studying
executable information fields as an alternative architectural direction for
long-horizon reasoning.

\begin{thebibliography}{99}

\bibitem{vaswani2017attention}
A. Vaswani et al., ``Attention Is All You Need,'' in \emph{Advances in Neural
Information Processing Systems}, 2017.

\bibitem{gu2022s4}
A. Gu, K. Goel, and C. R\'e, ``Efficiently Modeling Long Sequences with
Structured State Spaces,'' in \emph{Proceedings of ICLR}, 2022.

\bibitem{gu2023mamba}
A. Gu and T. Dao, ``Mamba: Linear-Time Sequence Modeling with Selective State
Spaces,'' in \emph{Proceedings of COLM}, 2024.

\bibitem{poli2023hyena}
M. Poli et al., ``Hyena Hierarchy: Towards Larger Convolutional Language
Models,'' in \emph{Proceedings of ICML}, pp. 28043--28078, 2023.

\bibitem{battaglia2018relational}
P. W. Battaglia et al., ``Relational inductive biases, deep learning, and graph
networks,'' \emph{arXiv:1806.01261}, 2018.

\bibitem{gilmer2017neural}
J. Gilmer et al., ``Neural Message Passing for Quantum Chemistry,'' in
\emph{Proceedings of ICML}, 2017.

\bibitem{bronstein2021geometric}
M. M. Bronstein et al., ``Geometric Deep Learning: Grids, Groups, Graphs,
Geodesics, and Gauges,'' \emph{arXiv:2104.13478}, 2021.

\bibitem{velickovic2022clrs}
P. Velickovic et al., ``The CLRS Algorithmic Reasoning Benchmark,'' in
\emph{Proceedings of ICML}, 2022.

\bibitem{bevilacqua2023causal}
B. Bevilacqua et al., ``Neural Algorithmic Reasoning with Causal
Regularisation,'' in \emph{Proceedings of ICML}, pp. 2272--2288, 2023.

\bibitem{zhu2021nbfnet}
Z. Zhu, Z. Zhang, L.-P. Xhonneux, and J. Tang, ``Neural Bellman--Ford Networks:
A General Graph Neural Network Framework for Link Prediction,'' in
\emph{Advances in Neural Information Processing Systems}, 2021.

\bibitem{kipf2018nri}
T. N. Kipf, E. Fetaya, K.-C. Wang, M. Welling, and R. S. Zemel, ``Neural
Relational Inference for Interacting Systems,'' in \emph{Proceedings of ICML},
pp. 2688--2697, 2018.

\bibitem{topping2022understanding}
J. Topping et al., ``Understanding Over-Squashing and Bottlenecks on Graphs via
Curvature,'' in \emph{Proceedings of ICLR}, 2022.

\bibitem{amari2016information}
S. Amari, \emph{Information Geometry and Its Applications}. Springer, 2016.

\bibitem{greydanus2019hamiltonian}
S. Greydanus, M. Dzamba, and J. Yosinski, ``Hamiltonian Neural Networks,'' in
\emph{Advances in Neural Information Processing Systems}, 2019.

\bibitem{finzi2020constraints}
M. Finzi, K. A. Wang, and A. G. Wilson, ``Simplifying Hamiltonian and
Lagrangian Neural Networks via Explicit Constraints,'' in \emph{Advances in
Neural Information Processing Systems}, 2020.

\bibitem{dijkstra1959note}
E. W. Dijkstra, ``A note on two problems in connexion with graphs,''
\emph{Numerische Mathematik}, vol. 1, pp. 269--271, 1959.

\bibitem{bellman1957dynamic}
R. Bellman, \emph{Dynamic Programming}. Princeton University Press, 1957.

\bibitem{chung1997spectral}
F. R. K. Chung, \emph{Spectral Graph Theory}. American Mathematical Society,
1997.

\bibitem{forman2003bochner}
R. Forman, ``Bochner's method for cell complexes and combinatorial Ricci
curvature,'' \emph{Discrete \& Computational Geometry}, vol. 29,
pp. 323--374, 2003.

\bibitem{efron1993bootstrap}
B. Efron and R. J. Tibshirani, \emph{An Introduction to the Bootstrap}.
Chapman \& Hall/CRC, 1993.

\bibitem{mcnemar1947note}
Q. McNemar, ``Note on the sampling error of the difference between correlated
proportions or percentages,'' \emph{Psychometrika}, vol. 12, pp. 153--157,
1947.

\end{thebibliography}

\appendix
\section{Operational Physics Glossary}

\begin{longtable}{@{}P{0.19\linewidth}P{0.36\linewidth}P{0.36\linewidth}@{}}
\caption{Reader-facing interpretation of physics-derived terms.}
\label{tab:glossary}\\
\toprule
Term & Implemented meaning & Scope \\
\midrule
\endfirsthead
\toprule
Term & Implemented meaning & Scope \\
\midrule
\endhead
$\Psi$ field & Token-indexed persistent neural state initialized from raw
semantics, position, and internal identity fibre. & Not a physical wavefunction
or literal field. \\
Density $\rho$ & Mean squared field amplitude per token. & Neural diagnostic,
not mass or energy density. \\
Adhesion $\Phi$ & Multiplicative gate over coherence, reconstruction,
stability, gradient, order, and boundary factors. & Constitutive modeling
choice, not a universal information law. \\
World lens & One of eight functions mapping the shared field to a directed edge
potential. & Functional multiworld dependence is tested; unique human semantic
roles are not claimed. \\
Metric/action & Non-negative learned path cost derived from world edge strength
and penalties. & Discrete path cost, not a proved smooth Riemannian metric. \\
Finite-speed cone & Hard admissibility condition based on path cost and a
discrete propagation interval. & Causal computational constraint, not physical
relativity. \\
GeoTransport & Geometry-conditioned movement of field and typed payload. & Not
standard QKV value mixing. \\
Product manifold & Discrete Cartesian state $(\text{token},\text{world})$. & A
product graph state space, not a proved differentiable manifold. \\
Geodesic & Minimum-action path computed by min-plus closure. & Shortest path in
the induced discrete action, not a continuum geodesic theorem. \\
Boundary deficit & Candidate-specific mismatch with registered terminal
conditions. & Hard decision observable, not a physical boundary measurement. \\
Boundary measurement & Final conjunction of typed register, terminal mass,
deficit, reachability, and action-margin tests. & Computational answer
acceptance or abstention rule. \\
Terminal mass & Transported support arriving at a valid terminal state. &
Payload support, not physical mass. \\
Memory geometry & Persistent path-dependent edge and payload state. & Bounded
task memory, not general episodic or language memory. \\
Curvature proxy & Detached Forman-style edge diagnostic. & Observer only; no
causal benefit claimed. \\
Spectrum & Eigenvalue diagnostics of a geometry-derived graph Laplacian. &
Observer only; no performance role claimed. \\
\bottomrule
\end{longtable}

\end{document}
